% !TEX root = ../main.tex
\section{Allokering av beslutningsrettigheter}\label{thr:decision-rights-allocation}

\paragraph{Definisjon.} \emph{Allokering av beslutningsrettigheter} er det første benet (\SOne) i \thref{organizational-architecture}{organizational architecture}: den eksplisitte tildelingen av myndighet til å initiere, ratifisere, gjennomføre og monitorere ulike beslutninger til personer, team og organer i virksomheten. Den grunnleggende designregelen er Jensen \& Mecklings (1992) \emph{co-location of knowledge and decision rights}: plasser beslutningsretten der den \thref{specific-vs-general-knowledge}{spesifikke videnen} er, så fremt agentkostnader og koordinasjonskostnader tillater det. [SLIDE/BOK]

\subsection*{Forklaring}

Alle beslutninger i en organisasjon må allokeres. Trade-off-en mellom to ekstremer driver designet:
\begin{itemize}
\item \textbf{Sentralisere} -- flytt informasjonen oppover til en sentral beslutningstaker. Fordeler: koordinering, skalafordeler, internalisering av eksternaliteter, mindre agentkostnad. Ulemper: dyrt å overføre spesifik viden, langsommere, demotiverende lokalt.
\item \textbf{Desentralisere} -- flytt beslutningen ned til der videnen er. Fordeler: bruker lokal viden, raskere, frigjør toppledelsens tid, motiverer. Ulemper: agentkostnader (lokal leder $\neq$ residual claimant), \thref{externalities}{eksternaliteter} mellom enheter, koordinasjonskostnader.
\end{itemize}
Optimal allokering ligger derfor mellom ytterpunktene og avhenger av: videnstruktur, agentkostnader, eksternaliteter, koordinasjonsbehov, og hvor godt \STwo\ og \SThree\ kan erstatte tapt direkte kontroll. [SLIDE]

\textbf{Det formelle valget} (BSZ kap.\ 12) brytes ned i fire spørsmål:
\begin{enumerate}
\item \emph{Hvilke} beslutninger skal allokeres? (pris, mengde, ansettelse, investering, strategisk valg).
\item \emph{Til hvilket nivå} (CEO, divisjonsleder, team, individ)?
\item \emph{Skal Management og Control separeres}? (jf.\ \thref{decision-management-control}{Decision Management vs.\ Decision Control}.)
\item \emph{Hvilken centertype} matcher allokeringen? (cost / revenue / profit / investment center; \thref{responsibility-centers}{responsibility centers}.)
\end{enumerate}

\textbf{Empowerment-vinklingen:} Slidens \emph{empowerment} er allokering nedover, men kun når (a) den ansatte har relevant spesifik viden, (b) det finnes måltall til å holde henne ansvarlig (S2), og (c) insentivene er innrettet (S3). Empowerment uten S2/S3 er meningsløst. [SLIDE]

\textbf{Optimal desentraliseringsgrad} (modell fra slide 11-12):
$NB = BD - AD - CD^2 \Rightarrow D^* = (B-A)/2C$, der $B$ = marginalnytte av desentralisering (lokal viden), $A$ = lineær kostnad (agentkostnad), $C$ = konveks kostnad (koordinasjon/eksternaliteter). Sentralt poeng: jo større lokal viden, jo høyere $D^*$; jo større eksternaliteter, jo lavere. [SLIDE]

\subsection*{Kjerneforutsetninger}
\begin{itemize}
\item Beslutninger kan identifiseres, separeres og tildeles eksplisitt.
\item Spesifik vs.\ generell viden er observerbar og stabil nok til design.
\item Måltall finnes (eller kan utvikles) som erstatning for direkte observasjon når man desentraliserer.
\item Insentiver kan kobles til måltallene (\SThree).
\item Eksternaliteter mellom enheter kan kartlegges og adresseres (priser, regler, sentralisering).
\item Organisasjonen kan endre allokeringen når omgivelsene endrer seg (ikke fanget i path dependence).
\end{itemize}

\subsection*{Muligheter}
\begin{itemize}
\item Gir et eksplisitt rammeverk for å diagnostisere organisasjonsproblemer: hvor sitter videnen, hvor sitter beslutningsretten, og samsvarer de?
\item Forklarer hvorfor IT/AI endrer optimal allokering: gjør tidligere spesifik viden generell $\Rightarrow$ rom for re-sentralisering (centralized dashboards, algorithmic pricing).
\item Forklarer franchise, divisjonalisering, og M-form: alle er svar på hvor beslutningsretten bør ligge.
\item Forklarer hvorfor noen beslutninger \emph{må} være sentrale (M\&A, kapitalstruktur, merkevare) -- store eksternaliteter / strategiske valg.
\item Begrunner separasjon av Management og Control: når desentraliserer man til ikke-residualkravshaver, må man legge til \thref{decision-management-control}{Control}-mekanismer.
\item Anvendelig på alle nivå: konsern $\to$ divisjon, divisjon $\to$ team, team $\to$ individ.
\end{itemize}

\subsection*{Begrensninger og kritikk}
\begin{itemize}
\item \textbf{Underspesifisert:} Modellen sier ``plasser beslutningsretten der videnen er'', men spesifikk grenseflate (hvilken eksakt beslutning) krever skjønn.
\item \textbf{Endogen viden:} Hvor videnen sitter, avhenger av tidligere allokering -- selvforsterkende. Sentraliserte selskaper utvikler ikke lokal viden.
\item \textbf{Strategisk:} Lokale ledere overdriver hvor spesifikk videnen er for å beholde makt. Sentralledere overdriver eksternaliteter for å re-sentralisere.
\item \textbf{Politiske kostnader:} \emph{Endring} av allokering møter motstand, lobbying og influence costs (Milgrom \& Roberts).
\item \textbf{Multi-task / interdependence:} Når oppgaver er vevd inn i hverandre, kan ikke beslutninger separeres rent -- modellen forutsetter klar dekomponering.
\item \textbf{Kvantifisering av modellen} ($D^* = (B-A)/2C$) er pedagogisk, men i praksis kan man sjelden måle $B$, $A$, $C$.
\item \textbf{Ignorerer kultur/identitet:} Stark organisasjonskultur kan erstatte formell allokering -- modellen er ren formell-rational.
\end{itemize}

\subsection*{Modell / illustrasjon}

\begin{figure}[h]
  \centering
  \includegraphics[width=0.65\linewidth]{BSZ_kap_12_2026/by_slide/slide_012/BSZ_kap_12_2026_slide_012_asset_01_image3.png}
  \caption{Trade-off mellom benefits og costs av desentralisering. Optimal $D^* = (B-A)/2C$. (BSZ kap.\ 12, slide 12) [SLIDE]}
\end{figure}

\begin{center}
\renewcommand{\arraystretch}{1.15}
\begin{tabular}{@{}p{3.7cm}p{4.4cm}p{5.4cm}@{}}
\toprule
\textbf{Driver} & \textbf{Trekker mot} & \textbf{Eksempel} \\
\midrule
Spesifik viden lokalt & Desentralisering & Lokal butikkpris (AutoMart) \\
Eksternaliteter mellom enheter & Sentralisering & Merkevare, IT, M\&A \\
Skalafordeler & Sentralisering & Innkjøp, treasury \\
Agentkostnader & Sentralisering & Når S2/S3 svake \\
Hastighet / lokal tilpasning & Desentralisering & Markeder med rask endring \\
Residual claimant lokalt & Desentralisering & Franchise \\
\bottomrule
\end{tabular}
\end{center}

\subsection*{Eksempler}
\begin{itemize}
\item \textbf{AutoMart (slide-eksempel):} Hvem skal sette pris -- CEO, regionsjef eller butikksjef? Lokal viden om kundepreferanser og konkurrenter ligger lokalt $\Rightarrow$ desentraliser til butikksjef, men sentraliser merkevaren og prislister-ramme. [SLIDE]
\item \textbf{Equinor E\&P:} Letevirksomhet desentralisert til prosjektledere (geologisk spesifikk viden), men store CAPEX-investeringer ratifiseres av styret (eksternalitet mot konsernrisiko, kapitalmarkedssignal). [EKS]
\item \textbf{Norsk Hydro -- pris vs.\ M\&A:} Pris på aluminium settes ikke lokalt (LME-marked), men teknologivalg på det enkelte smelteverk er desentralisert til verksledelse. M\&A-beslutninger ligger sentralt -- eksternaliteter for hele konsernet. [EKS]
\item \textbf{IKEA -- franchising vs.\ design:} Salgsbeslutninger desentralisert til varehus / franchisetaker (lokal viden). Design og produktkatalog er sentralisert i Älmhult (kodifiserbar viden + merkevareeksternalitet). [EKS]
\item \textbf{Spotify -- ``Squad model'':} Eksperiment med radikal desentralisering til autonome team. Senere delvis re-sentralisert -- viste at koordinasjonskostnader og eksternaliteter ble undervurdert. [EKS]
\item \textbf{Mærsk PtX-divisjonen:} Teknologi- og partnervalg desentralisert (spesifikk kjemisk viden), pricing av grønt drivstoff har store eksternaliteter mot Ocean-divisjonen $\Rightarrow$ koordinert; CAPEX over terskel ratifiseres sentralt. Klassisk balanseskap. [CASE]
\end{itemize}

\subsection*{Kobling til søyle og andre teorier}
Søyle: \SOne\ definisjonelt -- dette \emph{er} S1. Driver \STwo\ (måltall som speiler den allokerte beslutningen, jf.\ \thref{controllability-principle}{controllability}) og \SThree\ (insentiver på det som er allokert). Relaterte: \thref{specific-vs-general-knowledge}{Spesifik vs.\ generell viden} (kjernedriver), \thref{centralization-vs-decentralization}{Centralisering vs.\ desentralisering} (graden), \thref{decision-management-control}{Decision Management vs.\ Decision Control} (separasjon innen S1), \thref{externalities}{Eksternaliteter} (begrenser desentralisering), \thref{residual-claimants}{Residual claimants} (sentral for Control-allokering), \thref{agency-costs}{Agentomkostninger}, \thref{principal-agent}{Principal-agent}, \thref{moral-hazard}{Moral hazard}, \thref{responsibility-centers}{Responsibility centers} (implementering), \thref{organizational-architecture}{Organizational architecture} (helheten).

\paragraph{Brukt i spørsmål:} Q07.
